\section{Experiments}
\label{sec:method:experiment}

\subsection{Research Design}
\label{sec:method:experiment:design}
This is a controlled comparison against a fixed reference point. The manipulated variable is
the reduction method applied to the training data; the encoder, dataset, splits, label,
optimizer and evaluation protocol are held constant, so that differences in outcome are
attributable to the reduction. The reference point is the full-graph model of RQ1, and every
reduced configuration is reported relative to it rather than in isolation.

Two hypotheses are stated in advance (\ref{sec:intro:rq3},~\ref{sec:intro:rq4}): H1, that
path-contracting summarization retains more than edge-removing sparsification at matched
compression, because contraction shortens propagation paths where deletion only lengthens
them; and H2, that the domain-informed adaptations tested here are too simple to win
outright, in which case the informative outcome is what the null result says about where the
predictive signal lives.

\paragraph{Positive control.}
The design includes one configuration whose outcome is known in advance. Colour refinement at
count-cap $\infty$ on the exact-compression track is provably lossless for this encoder
(\ref{sec:method:architecture:exact}), so a model trained on those coarsened graphs must
score essentially identically to the full-graph baseline, both on reduced graphs and when
queried on full ones. This is what makes a poor cross-state result interpretable: without it,
RQ5 cannot distinguish ``models trained on reduced graphs do not generalize'' from ``the
evaluation path is wrong''. Any deviation of this configuration from the baseline is a defect
to be found, not a finding to be reported.

\paragraph{The experiments.}
\label{sec:method:experiment:phases}
Four experiments implement the design, and they do not map one-to-one onto the five research
questions (\ref{tab:experiment_phases}). RQ4 is a re-analysis of the Phase~3 runs under a
paired comparison rather than a separate experiment, which is why it costs no additional
training; RQ5 reuses the Phase~3 checkpoints and adds only an evaluation pass. The ordering
is a dependency chain rather than a preference: Phase~3 cannot start until Phase~2 has
produced the cached reduction artifacts, and Phase~4 needs Phase~3's checkpoints.

\begin{table}[htbp]
\centering
\caption{The four experimental phases, what each answers, and what each costs. Only
Phases~1 and~3 consume training compute.}
\label{tab:experiment_phases}
\small
\begin{tabular}{llll}
\toprule
\textbf{Phase} & \textbf{Answers} & \textbf{What is run} & \textbf{Training cost} \\
\midrule
1 & RQ1      & Full-graph model and baselines        & one run per model \\
2 & RQ2      & Offline reduction precompute, profiled & none \\
3 & RQ3, RQ4 & One model per reduced configuration    & one run per configuration \\
4 & RQ5      & Phase~3 checkpoints on full graphs     & none (evaluation only) \\
\bottomrule
\end{tabular}
\end{table}

Phase~1 establishes both the predictive reference point and the baseline hardware cost
(peak memory, epoch time) that every reduced configuration is measured against. Phase~2
requires no training and determines which methods are viable at corpus scale at all.
Phase~3 places the reduced configurations against the Phase~1 reference on a Pareto front
of predictive degradation versus memory and speed. Phase~4 queries models trained
exclusively on reduced graphs with unreduced ones, testing structural generalization.

\paragraph{Limits of the design.}
The comparison is within one encoder family, one synthesis algorithm and one bounded graph
scale. It can rank reduction methods against each other under those conditions; it cannot
establish that the ranking transfers to another architecture, another target algorithm, or
industrial-scale circuits. Section~\ref{sec:discussion:limitations} treats each in turn.

\subsection{Setup}
\label{sec:method:experiment:setup}

\subsubsection{Data Splitting}
\label{sec:method:experiment:splitting}
An 80/10/10 split by volume, partitioned \emph{by design}: all graphs derived from a given
source circuit fall in the same split. This is the strong choice and it is deliberate.
Because the corpus is built by transforming and then repeatedly optimizing a small number of
source designs, a random split would place a base graph in training and its own optimized
descendant in test: structurally near-identical circuits in different optimization states.
The model could then score well by recognising the design rather than by reading its
structure. Splitting by design removes that path entirely, at the cost of making the task
harder: the test set contains circuits the model has never seen in any state.

Two alternative strategies are implemented and available (split by optimization recipe, and
a plain random split). They are not used for the reported results; they are run once, at the
primary configuration, as the protocol-sensitivity experiment RQ1a
(\ref{sec:results:rq1:protocol}).
The three protocols differ in what they hold out, and each corresponds to an evaluation
setting already established in the literature (Table~\ref{tab:split_protocols}).

\begin{table}[htbp]
\centering
\caption{The three train/test protocols implemented in this work, ordered from leakiest to
strictest. Only the design-disjoint protocol is used for the reported results; the other two
are run once each at the same configuration to quantify how much of a reported score is
design recognition rather than structural reading.}
\label{tab:split_protocols}
\small
\begin{tabular}{@{}p{0.15\linewidth}p{0.30\linewidth}p{0.27\linewidth}p{0.12\linewidth}@{}}
\toprule
\textbf{Protocol} & \textbf{Held out} & \textbf{Established as} & \textbf{Role here} \\
\midrule
Random (per-row)
& Nothing. Individual graphs are assigned independently, so a base graph and its own
  optimized descendant may straddle the split.
& The common default in circuit representation-learning evaluations.
& Ablation \\[4pt]
Recipe-disjoint
& Whole synthesis recipes: every tier-0 step of a held-out recipe, plus its tier-1 and
  tier-2 descendants, across all designs. Circuits stay seen; the recipe does not.
& OpenABC-D Variant~1 \citep{chowdhury2021openabc}; LOSTIN transductive
  \citep{wu2022lostin}.
& RQ1a \\[4pt]
Design-disjoint
& Whole source circuits: every graph derived from a held-out design, in any optimization
  state, is confined to one split.
& OpenABC-D Variant~2 \citep{chowdhury2021openabc}; LOSTIN inductive
  \citep{wu2022lostin}.
& \textbf{Headline} \\
\bottomrule
\end{tabular}
\end{table}

% TODO: the random-split and recipe-split rows are placeholders until those runs land.

\subsubsection{Graph Scale Bounds}
\label{sec:method:experiment:scale}
The corpus is bounded at 366,040 nodes per AIG\@. This intermediate scale is what makes RQ1
possible: a single unreduced graph fits in accelerator memory, so the full-graph baseline
every other result is measured against can actually be trained. Larger circuits would remove
the reference point and leave only reduced configurations to compare with each other.

The bound is a scoping decision with a cost, and it is the limitation closest to the
motivation: the argument for this work invokes circuits of millions to billions of nodes, and
the evaluation does not reach them
(\ref{sec:discussion:limitations:scalability}).

\subsubsection{Dynamic Batching}
\label{sec:method:experiment:batching}
Graphs in this corpus vary in size by orders of magnitude, so a fixed batch size in graphs
gives wildly varying memory use and eventually fails on the largest circuits. Batches are
instead assembled to a maximum \emph{node} budget by greedy bin-packing, which bounds
activation memory regardless of which graphs are drawn. The batch plan is computed once and
cached, since replanning over a corpus this size is expensive.

Two budgets are used and they are deliberately distinct. Training uses the smaller of the
two; evaluation, which stores no activations or gradients, uses a substantially larger one.
Keeping them separate rather than raising a single shared value avoids invalidating the batch
plan the trained checkpoints were produced under.

\paragraph{Constraint on efficiency measurements.}
The evaluation budget must be identical across all configurations, because peak memory and
throughput both depend on it. Any change requires re-running every configuration; a mixed set
of budgets makes the RQ2 comparison meaningless rather than merely noisy. This is stated here
because it constrains how the results in~\ref{sec:results:rq2} may be read and combined.

One irreducible peak remains that no budget can lower: a graph larger than the budget still
forms a batch of its own, since a graph-level target cannot be split across batches. The
largest circuit in the corpus therefore sets a floor on peak memory for every configuration
that does not reduce node counts.

\subsubsection{Model Configuration}
\label{sec:method:experiment:model}
An Edge-Aware Graph Convolutional Network (GCN+) is trained to predict the optimizability of
a circuit when subjected to the \texttt{Orchestrate} synthesis script. The architecture is
described in~\ref{sec:method:architecture} and is held fixed across all reduction
configurations; the baselines it is compared against are described
in~\ref{sec:method:architecture:baselines}.

\subsubsection{Hyperparameter Tuning}
\label{sec:method:experiment:hp}
Hyperparameters were selected by an Optuna \citep{akiba2019optuna} sweep on the held-out
balanced subset of~\ref{sec:method:data:hpsubset}, which is then excluded from all
subsequent training and evaluation. Tuning on a subset that later appeared in training or
test would let the reported scores absorb information from the tuning process; excluding it
entirely costs 50,000 graphs and removes the question.

Tuning was performed once, on the full-graph configuration, and the resulting values are held
fixed for every reduced configuration. This is the conservative choice for the comparison
being made: re-tuning per reduction method would confound ``this reduction retains more
signal'' with ``this reduction received more tuning effort''. It does mean reduced
configurations are evaluated at hyperparameters chosen for a different input distribution,
which if anything biases against them. That is worth stating explicitly, since it bounds the
strength of any negative result about reduction.
% TODO: record the search space and the number of trials from src/shell/big_hp_tuning.sh.

\subsubsection{Training Configuration}
\label{sec:method:experiment:training}
Smooth L1 loss with $\beta = 0.01$, chosen well below the default because the target occupies
a narrow band within $[0,1]$ and the default threshold would place nearly every residual in
the quadratic regime, making the loss effectively squared error. Optimization uses AdamW with
learning rate $3 \times 10^{-4}$ and weight decay $3.96 \times 10^{-3}$, a linear warmup over
the first 10,000 steps, and \texttt{ReduceLROnPlateau} thereafter (factor 0.5, patience 2,
floor $10^{-6}$). Training stops early on validation loss with patience 3. Gradients are
clipped at 1.0.

Two settings affect the measurements rather than the model and must therefore be disclosed
with them: \texttt{torch.compile} is enabled, which changes every timing figure, and mixed
precision is used on GPU, which changes every memory figure. Neither is varied across
configurations, so comparisons remain valid; absolute numbers are only meaningful together
with these settings.

\subsubsection{Hardware and Software Environment}
\label{sec:method:experiment:hardware}
All experiments ran on the Snellius national supercomputer: H100 80\,GB GPU nodes for
training and evaluation, and CPU-only nodes for the offline reduction precompute, which is
embarrassingly parallel over graphs and does not benefit from a GPU\@.

Two measured characteristics of the workload are worth recording, because they explain the
shape of the RQ2 results rather than merely describing the setup. First, memory scales close
to linearly with nodes and edges per batch, the encoder's own parameters being negligible at
this width and depth, which is why a node budget bounds memory effectively at all. Second,
message passing is latency-bound on gather and scatter operations rather than
compute-bound: at full device occupancy the floating-point pipelines and memory bandwidth
both sit far below saturation. The consequence for RQ2 is that a larger
batch amortises per-batch overhead but does not deliver a proportional speedup, and that
reductions which remove \emph{edges} should help throughput more than their node compression
alone would suggest.

% TODO: library versions, the synthesis tool version, and the total wall-clock budget.
